\chapter{State of the Art}
\label{chap:state-of-the-art}

\section{Clustering Foundations}
\label{sec:clustering-foundations}

\subsection{K-Means Clustering}
k-means is a partition-based algorithm that divides a dataset $X = \{x_1, x_2, ..., x_N\}$ into $k$ non-overlapping clusters by minimizing the within-cluster sum of squares \cite{macqueen1967}. Let each observation satisfy $x_i \in \mathbb{R}^d$, where $d$ denotes the number of variables or features. The objective function is:
\begin{equation}
  J = \sum_{j=1}^{k} \sum_{x_i \in S_j} ||x_i - \mu_j||^2
\end{equation}
where $\mu_j$ is the centroid of cluster $S_j$.

Its standard iterative procedure alternates between assigning each point to the nearest centroid and recomputing centroids from the assigned points. k-means is attractive because of its computational efficiency, but it assumes clusters are approximately spherical, compact, and separable in Euclidean space. It is also sensitive to initialization and to outliers.

\subsection{Agglomerative Hierarchical Clustering}
Hierarchical Clustering constructs a nested family of clusterings represented by a dendrogram. This study uses the agglomerative variant, where each observation begins as a singleton cluster and clusters are merged iteratively according to a linkage rule \cite{ward1963}.

The linkage rule adopted here is \textbf{Ward's method}, which chooses the merge that minimizes the increase in within-cluster variance. For clusters $A$ and $B$, the merge cost is:
\begin{equation}
  D(A, B) = \text{ESS}_{A \cup B} - (\text{ESS}_A + \text{ESS}_B)
\end{equation}
where $\text{ESS}_X = \sum_{x \in X} ||x - \mu_X||^2$.

Compared with k-means, Hierarchical Clustering is computationally more expensive, but it offers two relevant advantages for this report: it is deterministic and it can better preserve connectivity-based structure in non-convex or skewed settings.

\subsection{Comparative Summary}
Table~\ref{tab:kmeans_vs_hc} summarizes the main operational differences between the two clustering strategies considered in this study.

\begin{table}[H]
  \centering
  \caption{Comparison of K-Means and Hierarchical Clustering}
  \label{tab:kmeans_vs_hc}
  \renewcommand{\arraystretch}{1.2}
  \begin{tabular}{@{}lp{5.8cm}p{5.8cm}@{}}
    \toprule
    \textbf{Feature} & \textbf{k-means} & \textbf{Hierarchical (Ward)} \\ \midrule
    Complexity & Approximately linear in $N$ for fixed $k$ and iteration budget & Between $O(N^2)$ and $O(N^3)$ depending on implementation \\
    Geometry & Assumes spherical and convex clusters & More flexible, based on connectivity and variance increase \\
    Parameters & Requires $k$ in advance & Allows a post-hoc choice of $k$ from the dendrogram \\
    Robustness & Sensitive to initialization, outliers, and skew & More robust to irregular shapes, though computationally heavier \\
    Deterministic & No & Yes \\ \bottomrule
  \end{tabular}
  \source{Compiled by the author}
\end{table}

\section{Synthetic Data Utility}
\label{sec:synthetic-data-utility}

Synthetic data has become a common strategy for reconciling privacy constraints with analytical demand. The promise is straightforward: replace restricted real records with generated observations that preserve enough statistical structure to support downstream analysis without exposing personal information \cite{nowok2016,qian2024}.

For clustering, however, utility is not just a matter of matching marginal distributions. The generated data must preserve cluster geometry, balance, separability, and dispersion. A synthetic dataset can look statistically plausible overall while still failing to support valid unsupervised analysis if cluster boundaries are distorted.

\subsection{CART-Based Synthesis}
The synthetic datasets in this report are generated with the CART method from the \texttt{synthpop} package \cite{nowok2016}. CART synthesis models variables sequentially through conditional regression trees. This gives it flexibility for non-linear relationships, but also introduces a characteristic artifact: smooth boundaries can become orthogonal and block-like because decision trees partition space through axis-aligned splits.

That property is especially relevant when the downstream task depends on distance geometry. Centroid-based algorithms may react poorly to staircase-like boundaries, while connectivity-based methods may remain more stable. This distinction is also why the present work separates the fidelity of the synthetic data from the robustness of the clustering algorithm itself.

\section{Related Work and Positioning}
\label{sec:related-work}

This report extends previous synthetic-data utility work by moving from general utility claims to a clustering-specific analysis. Prior evaluations often emphasize supervised performance or broad descriptive similarity. Those criteria do not directly answer whether unsupervised methods still recover meaningful structure after synthesis.

The present study therefore positions utility as an interaction between:
\begin{itemize}
  \item the \textbf{generation method}, which can preserve or distort geometry; and
  \item the \textbf{analysis method}, which may be more or less sensitive to those distortions.
\end{itemize}

This framing also differentiates the present report from published clustering benchmarks such as MDCGen \cite{iglesias2019}, which focus on generating synthetic cluster scenarios for algorithm testing. Here the question is different: starting from a real reference dataset, does privacy-motivated synthesis preserve enough structure for downstream unsupervised analysis? Rather than asking whether synthetic data is useful in general, this study asks \textit{for which clustering algorithms, under which distributional conditions, and according to which structural criteria} synthetic data remains useful.
